\subsection{p 进李群}

这样的群最早于 1907 年出现在亨塞尔关于 p 进解析函数（由整级数展开定义）的工作{[}157 e{]}中。他特别研究了指数函数与对数函数；尽管定义它们的级数的行为先验地看来令人惊讶（例如指数级数并非处处收敛），它们的函数性质仍然成立，这就给出了 \(Q_{p}\) 的加法群与乘法群之间的一个局部同构（或者更一般地，给出特征为零的每个完备超度量域的加法群与乘法群之间的局部同构）。

A. 韦伊{[}330 a{]}与 E. 卢茨{[}210{]}关于 p 进椭圆曲线（1936 年）的工作，所关心的同样是交换群（不过这一次不是线性的）。除算术应用之外，其中还可以找到该群与加法群之间一个局部同构的构造，它建立在一个不变微分形式的积分之上。这一方法同样可以应用于阿贝尔簇，C. 沙博蒂不久之后指出了这一点，他未作更多说明便使用这一方法证明了“莫德尔猜想”的一个特例{[}61{]}。

从这时起人们已经清楚，李群的局部理论几乎可以不加改变地应用于 p 进情形。李群—李代数“词典”的基本定理是 1942 年在 R. 胡克（谢瓦莱的学生）的博士论文{[}166{]}中建立的；这项工作还包含关于实李群闭子群的 E. 嘉当定理的 p 进类似物。

更近些时候，M. 拉扎尔{[}195 b{]}对 \(\mathbf{Q}_p\) 上的紧解析群发展了“词典”的一种更精确的形式。他证明，\(p\) 进解析结构在紧群 \(G\) 上的存在与 \(G\) 上某些过滤的存在紧密相连，并给出了它的多种应用（例如对 \(G\) 的上同调）。拉扎尔的工具之一是对邓金关于 \(p\) 进豪斯多夫级数收敛性结果的一个改进{[}195 a{]}。

